% Generated draft: requires mathematical and source review.
\documentclass[11pt]{article}
\usepackage[margin=25mm]{geometry}
\usepackage{fontspec}
\setmainfont{Latin Modern Roman}
\usepackage{amsmath,amssymb,mathtools,unicode-math}
\setmathfont{Latin Modern Math}
\usepackage{enumitem,needspace,xurl,hyperref}
\setlength{\parindent}{0pt}
\setlength{\parskip}{5pt}
\setlength{\emergencystretch}{3em}
\newcommand{\N}{\mathbb{N}}
\newcommand{\Z}{\mathbb{Z}}
\newcommand{\Q}{\mathbb{Q}}
\newcommand{\R}{\mathbb{R}}
\newcommand{\C}{\mathbb{C}}
\newcommand{\F}{\mathbb{F}}
\newcommand{\A}{\mathbb{A}}
\newcommand{\PP}{\mathbb{P}}
\newcommand{\E}{\mathbb{E}}
\newcommand{\eps}{\varepsilon}
\newcommand{\dd}{\,\mathrm{d}}
\begin{document}
\section*{UnsolvedMath ID: 800034}
\textbf{Source problem code:} catalogue entry
\subsection{Definitions and mathematical statement}
A semigroup is a nonempty set with an associative binary operation. A semigroup \(S\) is commutative if \(ab=ba\) for all \(a,b\in S\), and periodic if
\[
\forall s\in S\;\exists m,n\in\N\quad 1\le m<n\ \text{ and }\ s^m=s^n.
\]
An Abelian group \(G\), written multiplicatively with identity \(1_G\), is periodic if
\[
\forall g\in G\;\exists n\in\N\quad n\ge1\ \text{ and }\ g^n=1_G.
\]
For subsets \(A,B\subseteq G\), define their setwise product by
\[
AB=\{ab:a\in A,\ b\in B\}.
\]
The source does not specify whether its power semigroup includes the empty subset. To display the two possible conventions, put
\[
\mathcal P_0(G)=\{A:A\subseteq G\},\qquad
\mathcal P_1(G)=\{A\subseteq G:A\ne\varnothing\},
\]
each equipped with setwise multiplication.

For a fixed choice \(\delta\in\{0,1\}\), the corresponding formulation of Problem 36 asks: for every periodic commutative semigroup \(S\), do there exist a periodic Abelian group \(G\) and an injective map
\[
\varphi:S\longrightarrow\mathcal P_\delta(G)
\]
such that
\[
\forall a,b\in S,\qquad \varphi(ab)=\varphi(a)\varphi(b)?
\]
The two values of \(\delta\) represent alternative readings of the source, which does not select one explicitly.

For the supplied background qualification, a semigroup is regular if
\[
\forall a\in S\;\exists x\in S,\qquad a=axa.
\]
The record reports that the assertion is true for all regular commutative semigroups and for some finite non-regular commutative semigroups, citing \([\mathrm{Bersh+Kubl}]\). This is a supplied background claim, not an independently verified result. Its wording does not explicitly repeat the periodicity hypothesis for the regular case; the intended scope and the missing bibliographic reference require review.
\subsection{Short English statement}
Does every periodic commutative semigroup admit an injective homomorphism into a semigroup of subsets of some periodic Abelian group, with multiplication given by setwise products? The source leaves the empty-subset convention unspecified. It labels the problem open and reports affirmative regular and selected finite non-regular cases; these supplied assessments require review.
\subsection{Sources}
\begin{itemize}
\item Local UnsolvedMath registry, supplied complete record: permanent ID \(800034\), exact source problem code \texttt{catalogue entry}, title \emph{The Lyapin's Notebook --- Problem 36}. Upstream revision: \texttt{d371ee917ac9647fbb06cd6fc529deec6f3dda82}. The registry labels the problem \texttt{open}; this label is supplied metadata, not independent verification.
\item Bershadsky, Kublanovsky, Mashevitzky, \emph{The Lyapin's Notebook: Unsolved Problems in Semigroup Theory} (2026), Problem 36. Supplied source URL: \url{https://arxiv.org/abs/2604.04763}. The provenance records access on 2026-10-16 and extraction by \texttt{arxiv\_latex}; these are supplied provenance statements, not a claim of access during this preparation.
\item The problem statement attributes the question specifically to S.G. Bershadsky and S.I. Kublanovsky. The registry's \texttt{proposed\_by} field lists Bershadsky, Kublanovsky, Mashevitzky, with proposed year 2026; these two attribution scopes are retained separately.
\item The source cites \([\mathrm{Bersh+Kubl}]\) for its reported affirmative cases. No bibliographic expansion or separate link for that reference is supplied, so the underlying result cannot be identified precisely from this record alone.
\end{itemize}
\end{document}
